\documentclass[11pt]{article}

\usepackage[margin=1in]{geometry}
\usepackage{amsmath,amssymb,amsthm}
\usepackage[expansion=false]{microtype}
\usepackage{xcolor}
\usepackage{hyperref}

\hypersetup{
  colorlinks=true,
  linkcolor=blue!50!black,
  citecolor=blue!50!black,
  urlcolor=blue!50!black
}

\newtheorem{theorem}{Theorem}
\newtheorem{lemma}[theorem]{Lemma}
\theoremstyle{definition}
\newtheorem{definition}[theorem]{Definition}
\newtheorem*{remark}{Remark}

\newcommand{\R}{\mathbb{R}}
\newcommand{\Z}{\mathbb{Z}}
\newcommand{\Tr}{\mathrm{Tr}}

\title{The Subset-Vanishing Theorem: $B \equiv 0 \;\Rightarrow\; B_i \equiv 0$\thanks{Section draft, intended as \S3 of the rewritten \emph{Locality and Unitarity of $\Tr(\phi^3)$ Tree Amplitudes from Hidden Zeros}.}}
\author{Jony}
\date{April 2026}

\begin{document}
\maketitle

\setcounter{section}{3}
\renewcommand{\thesection}{\arabic{section}}
\renewcommand{\thesubsection}{\thesection.\arabic{subsection}}

\section*{Section 3 --- The Subset-Vanishing Theorem}

\noindent
This section gives a one-paragraph proof of the implication
\[
B(x,y) \equiv 0 \text{ on } \{\ell(x) = 0\} \quad\Longrightarrow\quad B_i(x,y) \equiv 0 \text{ on } \{\ell(x) = 0\} \text{ for every } i,
\]
where $B$ is partitioned by homogeneity in the $x$-variables and $\ell$ is any linear map of $x$ alone. This is the result that lets one work zone-by-zone in the BCFW grading without losing information, and it is the workhorse of every uniqueness argument in the bootstrap. The reason for stating it as a stand-alone theorem with minimal hypotheses is that the standard reference \cite[eqs.~(15)--(18)]{Rodina} reaches the same conclusion through a chain of structural claims about $X^{(0)}$ vs.\ $X^{(\infty)}$ identities (the latter requiring his Appendix~A); the proof below shows that this chain is unnecessary --- a single rescaling argument suffices.

\subsection{Statement and proof}

\begin{definition}[Homogeneity partition]
Let $B(x,y)$ be a rational function with $x \in \R^n$, $y \in \R^m$, and finitely many monomials. Group the monomials of $B$ by total degree in $x$:
\[
B(x,y) \;=\; \sum_{i \in \Z} B_i(x,y), \qquad B_i(\lambda x,\, y) \;=\; \lambda^i \, B_i(x,\, y) \text{ for all } \lambda \neq 0.
\]
The sum is finite. We call $\{B_i\}$ the \emph{$x$-homogeneity partition} of $B$.
\end{definition}

\begin{theorem}[Subset vanishing]\label{thm:subset}
Let $B(x,y)$ be a rational function with $x$-homogeneity partition $B = \sum_i B_i$, and let $\ell : \R^n \to \R^k$ be a linear map (so $\ell$ depends on $x$ only). Then
\[
\Bigl[\, \ell(x) = 0 \;\Rightarrow\; B(x,y) = 0 \,\Bigr]
\qquad\Longleftrightarrow\qquad
\Bigl[\, \ell(x) = 0 \;\Rightarrow\; B_i(x,y) = 0 \text{ for every } i\, \Bigr].
\]
\end{theorem}

\begin{proof}
``$\Leftarrow$'' is trivial since $B = \sum_i B_i$.

For ``$\Rightarrow$'', fix any $(x,y)$ with $\ell(x) = 0$ and any nonzero $\lambda \in \R$. By linearity, $\ell(\lambda x) = \lambda\,\ell(x) = 0$, so $(\lambda x, y)$ also lies on the zero locus. The hypothesis gives $B(\lambda x, y) = 0$. Expanding,
\[
0 \;=\; B(\lambda x, y) \;=\; \sum_{i} B_i(\lambda x, y) \;=\; \sum_{i} \lambda^{i}\, B_i(x, y).
\]
The right-hand side is a Laurent polynomial in $\lambda$ that vanishes at every nonzero $\lambda$. A Laurent polynomial vanishing on an infinite set has every coefficient equal to zero, so $B_i(x,y) = 0$ for every $i$. \qedhere
\end{proof}

\begin{remark}
The proof uses only that $\ell$ acts linearly on the homogeneity variables. Nothing is assumed about the structure of $\ell$, the relations the $x_i$ satisfy on the locus $\{\ell(x) = 0\}$, or the algebraic form of $B$. In particular, no claim is made or needed that $B$ lies in the ideal generated by $\ell$.
\end{remark}

\subsection{Connection to Rodina's eqs.~(15)--(18)}

In \cite[``Subset zeros and UV scaling'']{Rodina}, the same equivalence appears in the specific setting where $B$ is a rational function of planar Mandelstams $\{X_{ij}\}$, the locus is a hidden $k$-zero $\{c_{ij} = 0\}$, and the homogeneity grading is the BCFW-shift scaling under $p_{k+1} \to p_{k+1} + zq$, $p_n \to p_n - zq$. With the dictionary
\[
x \;=\; (\text{$z$-dependent } X_{ij}), \qquad y \;=\; (\text{$z$-independent } X_{ij}), \qquad \ell(x) \;=\; \text{the $c_{ij}$ relations restricted to $x$},
\]
Theorem~\ref{thm:subset} is precisely his eq.~(14), restricted to ``$B$ satisfies the zero $\Leftrightarrow$ each $B_i$ satisfies the zero.''

\paragraph{Rodina's argument.}
The chain in \cite[eqs.~(15)--(18)]{Rodina} runs through the structural form of $B$ on the zero locus. Eq.~(15) asserts that $B(X^{(0)}_{ij}) = 0$ implies $B = \sum_{c_{ij} \in \text{zero}} c_{ij}\,\{\dots\}_{ij}$, on the grounds that ``the only identities the $X^{(0)}_{ij}$ satisfy are given by $c_{ij}=0$.'' BCFW-shifting and taking $z \to \infty$ (eq.~16) gives $B \to z^m B_m(X^{(\infty)}_{ij})$. Eq.~(17) then claims that $B_m$ also has a representation $\sum c_{ij}\{\dots\}_{ij}$ \emph{with the same $c_{ij}$}, on the grounds that ``$X^{(\infty)}_{ij}$ and $X^{(0)}_{ij}$ satisfy identical relations'' --- a non-obvious claim established by the basis-counting in his Appendix~A. From eq.~(17) the leading $B_m$ peels off, the procedure iterates on $B - B_m$ (eq.~18), and one extracts each $B_i$ in turn.

The argument is correct, but it has three moving parts: (1) the structural representation in eq.~(15); (2) the identification of $X^{(\infty)}$ relations with $X^{(0)}$ relations (Appendix~A); (3) the iterative peel-off (eqs.~17--18).

\paragraph{What the rescaling proof eliminates.}
Theorem~\ref{thm:subset} reaches the same conclusion using only one move: rescale $x \mapsto \lambda x$ inside the zero locus (legal because $\ell$ is linear), then read off Laurent coefficients in $\lambda$. In particular:
\begin{itemize}
  \item No structural representation $B = \sum c_{ij}\{\dots\}_{ij}$ is invoked. The hypothesis is the bare equation $\ell(x) = 0 \Rightarrow B = 0$; we never write $B$ as an explicit element of any ideal.
  \item No comparison of $X^{(0)}$ and $X^{(\infty)}$ relations is required. The argument stays inside one fixed locus throughout. \emph{Appendix~A of \cite{Rodina} is therefore not on the critical path for this equivalence.}
  \item No iteration is needed. All $B_i$ vanish simultaneously from a single application of the Laurent argument.
\end{itemize}

\paragraph{Where Appendix~A is still needed.}
\cite[Appendix~A]{Rodina} proves that $X^{(\infty)}_{ij}$ and $X^{(0)}_{ij}$ satisfy the same set of $c_{ij}=0$ identities. This fact remains useful elsewhere in his program --- specifically for the third equivalence in eq.~(14), the statement that subsets exhibit \emph{enhanced} UV scaling $z^{i-1}$ rather than $z^i$. Theorem~\ref{thm:subset} does not reach the enhanced-scaling refinement; it gives only ``$B$ satisfies the zero $\Leftrightarrow$ each $B_i$ does.'' The enhanced-scaling step compares the $z^{m-1}$ coefficient on the limit locus to the $c_{ij}$ ideal on the zero locus, and that comparison genuinely needs the Appendix~A identification. We make no claim against that part of \cite{Rodina}; the point is to make the dependency explicit and isolated.

\begin{thebibliography}{9}
\bibitem{Rodina} L.~Rodina, ``Hidden zeros are equivalent to enhanced ultraviolet scaling and lead to unique amplitudes in $\Tr(\phi^3)$ theory,'' arXiv:2406.04234.
\end{thebibliography}

\end{document}
